\chapter{First Aid Kit Checklist}\label{ch:firstaid_kit}

This checklist provides guidance for assembling a comprehensive home/personal first aid kit. Adjust contents based on specific needs, activities, and family composition.

\begin{center}
\begin{tabular}{@{}lll@{}}
\hline
\textbf{Category} & \textbf{Item} & \textbf{Quantity} \\
\hline
Wound Care & Adhesive bandages (assorted sizes) & 20+ \\
& Sterile gauze pads (4x4 inches) & 20 \\
& Sterile adhesive tape & 1 roll \\
& Medical scissors (blunt end) & 1 \\
& Tweezers & 1 pair \\
& Antibiotic ointment packets & 5-10 \\
& Hydrocortisone cream (1\%) & 1 tube \\
& Antiseptic wipes (alcohol/povidone-iodine) & 10-20 \\
& Triple antibiotic ointment & 1 tube \\
& Burn ointment/gel & 1 tube \\
& Elastic compression wrap (ACE bandage) & 1-2 \\
& Triangular bandages & 2 \\
& Safety pins & Several \\
& Non-latex gloves (pair) & 2-4 pairs \\
& Instant cold pack & 1-2 \\
Thermometer & Digital thermometer & 1 \\
Medications & Pain relievers (acetaminophen/ibuprofen) & Supply as needed \\
& Antihistamine (diphenhydramine) & As needed \\
& Anti-diarrhea medication & As needed \\
& Laxative (optional) & As needed \\
& Personal prescription medications & Adequate supply \\
Tools & Flashlight with extra batteries & 1 \\
& Battery-powered or hand-crank radio & 1 \\
& Whistle (for signaling) & 1 \\
& Duct tape & 1 roll \\
& Plastic sheet/tarp & 1 \\
& Manual CPR face shield/barrier device & 1 \\
Documentation & Emergency contacts list & 1 copy \\
& First aid reference guide & 1 \\
& Insurance cards/Copies of IDs & Copies \\
& Medical history forms & As needed \\
\hline
\end{tabular}
\end{center}

\quicknote{\textbf{Special Considerations}: Include infant/child-specific items for families with children (fever reducers sized correctly, childproof lockouts, pediatric pain medications). Include pet supplies separately. Consider adding extra items for travel/vehicle kits (safety triangles, jumper cables, blanket). Store all medications in a cool, dry place away from direct sunlight. Check expiration dates every 6 months and replace as necessary.}

\section{How to Use Your First Aid Kit Items}

Knowing how to use each item correctly is as important as owning it.

\begin{description}
  \item[Non-latex gloves] Put them on before touching blood, wounds, or body fluids. Remove by peeling the cuff and turning them inside out; dispose of them in a sealed bag. Never reuse or wash gloves.
  \item[Antiseptic wipes] Use on intact skin around a wound to reduce skin bacteria before dressing. Do NOT pour into deep wounds -- cleaning the surrounding skin with clean water and mild soap is sufficient inside a wound.
  \item[Sterile gauze pads] Apply directly over a wound and hold with firm pressure to control bleeding. A blood-soaked dressing should be left in place and added to, never removed.
  \item[Adhesive tape] Secures gauze and dressings. Use strips long enough to anchor the dressing without wrapping so tightly that circulation is restricted.
  \item[Elastic (ACE) wrap] Provides compression for sprains and strains and holds dressings in place. Wrap snugly but not so tightly that it causes numbness, tingling, or color change below the wrap; unwrap and rewrap if it feels too tight.
  \item[Triangular bandages] Can be folded as an arm sling, folded into a broad bandage for pressure dressings, or rolled for padding. Improvised with any clean cloth if missing.
  \item[Medical scissors] Cut bandages, tape, and clothing around an injury. Blunt-tipped versions are safer near skin.
  \item[Tweezers] Remove splinters and ticks. Use fine-tipped tweezers for ticks, gripping close to the skin and pulling straight out steadily. Sterilize the tips with antiseptic before and after.
  \item[Instant cold pack] Activates when squeezed or shaken. Wrap in a cloth before placing on skin; apply for 15--20 minutes every 2--3 hours for sprains, strains, and bruises. Never apply directly to skin (frostbite risk).
  \item[Digital thermometer] Take temperature in the mouth, armpit, or ear per the device instructions. Note the reading and time for responders. Clean the probe after each use.
  \item[Antibiotic ointment] Apply a thin layer to clean minor cuts and scrapes before dressing. Do not use on deep, large, or infected wounds without medical advice.
  \item[Hydrocortisone cream (1\%)] For itchy insect bites and rashes. Apply a thin layer; do not use on broken skin or for more than a few days without medical advice.
  \item[Pain relievers] Acetaminophen and ibuprofen reduce pain and fever. Follow the label dosage, do not combine beyond instructions, and do not give aspirin to children or teenagers. Ask about allergies and other medications first.
  \item[Antihistamine] For allergic reactions and itching. Follow the label; some cause drowsiness (diphenhydramine) and affect driving. For severe allergic reactions, use the person's epinephrine auto-injector and call emergency services -- an antihistamine is not a substitute.
  \item[Manual CPR face shield / barrier device] Place over the person's mouth and nose before rescue breaths to reduce infection risk, as described in Chapter~\ref{ch:cpr}.
  \item[Flashlight and radio] Use to assess a scene and stay informed. Keep spare batteries sealed and replace them annually.
  \item[Whistle] Carry it on outdoor trips for signaling; three blasts is a recognized distress signal.
  \item[Duct tape and plastic sheet] Improvise waterproofing, splints, slings, and shelters in extended or wilderness situations.
\end{description}

\section{Medication Safety}

\begin{itemize}[leftmargin=*]
  \item Read and follow the label: dose, timing, and maximum daily amount.
  \item Check the expiration date before use and replace expired medication.
  \item Confirm allergies and ask whether the person has been advised against a medication by their doctor.
  \item Keep all medications in their original containers and out of reach of children.
  \item Note exactly what was given, when, and how much -- report this to responders.
\end{itemize}

\section{Kit Maintenance Checklist}

\begin{enumerate}[leftmargin=*]
\item Check expiration dates every 6 months; replace used or expired items immediately
\item Restock after every use -- never leave a kit incomplete
\item Rehearse with the kit: know where each item is and how to use it
\item Tailor the kit to your activities (home, vehicle, workplace, travel, outdoor)
\item Store the kit in an accessible, labeled place that everyone in the household knows
\item Revisit the kit after seasonal changes (add cold-weather items in winter, allergy items in spring)
\end{enumerate}
